\chapter{Hạ tầng phát triển và triển khai}

\section{Cấu trúc repo và công cụ phát triển}

Dự án dùng Python cho backend/NLP pipeline và Next.js cho giao diện web. Dependency Python được quản lý bằng \texttt{pyproject.toml}; version control dùng Git. Cấu trúc repo đáp ứng yêu cầu đồ án:

Bảng \ref{tab:project-structure} liệt kê các thư mục chính và vai trò của từng nhóm mã nguồn/tài nguyên trong repo.

\begin{table}[H]
\centering
\small
\caption{Cấu trúc dự án}
\label{tab:project-structure}
\begin{tabularx}{\textwidth}{@{}p{3.0cm}X@{}}
\toprule
\textbf{Đường dẫn} & \textbf{Mục đích} \\
\midrule
\texttt{src/} & Backend, NLP pipeline, monitoring và MLOps code. \\
\texttt{web/} & Next.js frontend cho upload, review và feedback. \\
\texttt{data/} & Dữ liệu đánh giá, dữ liệu huấn luyện hoặc đường dẫn dữ liệu runtime. \\
\texttt{models/} & Checkpoint baseline đã huấn luyện và output runtime được ignore. \\
\texttt{configs/} & Cấu hình YAML runtime, được load qua \texttt{CONFIG\_FILE}. \\
\texttt{tests/} & Unit test và smoke test cho API, config, guardrails, MLOps và baseline training. \\
\bottomrule
\end{tabularx}
\end{table}

\section{Kiến trúc runtime}

Hệ thống có thể triển khai bằng Docker Compose. Profile mặc định khởi động Next.js, FastAPI, Celery worker, Redis, PostgreSQL, Ollama, Prometheus và Grafana. Profile MLOps bổ sung Celery Beat, trainer worker và MLflow.

Hình \ref{fig:runtime} mô tả luồng runtime từ giao diện người dùng đến API, hàng đợi, worker xử lý nền và lớp monitoring. Bảng \ref{tab:runtime-services} mô tả trách nhiệm triển khai của từng service trong Docker Compose.

\begin{figure}[H]
\centering
\begin{tikzpicture}[
  node distance=0.45cm,
  procstep/.style={draw=primary, fill=light, rounded corners=3pt, text width=6.0cm, minimum height=0.58cm, align=center, font=\scriptsize},
  arrow/.style={-Latex, thick, color=darkgray}
]
  \node[procstep] (web) {Next.js Web UI};
  \node[procstep, below=of web] (api) {FastAPI API + PostgreSQL state};
  \node[procstep, below=of api] (queue) {Redis queue/cache};
  \node[procstep, below=of queue] (worker) {Celery worker: ASR, diarization, LLM, guardrails};
  \node[procstep, below=of worker] (obs) {Prometheus/Grafana monitoring};
  \draw[arrow] (web) -- (api);
  \draw[arrow] (api) -- (queue);
  \draw[arrow] (queue) -- (worker);
  \draw[arrow] (worker) -- (obs);
\end{tikzpicture}
\caption{Kiến trúc runtime}
\label{fig:runtime}
\end{figure}

\begin{table}[H]
\centering
\small
\caption{Các service triển khai trong Docker Compose}
\label{tab:runtime-services}
\begin{tabularx}{\textwidth}{@{}p{2.8cm}X@{}}
\toprule
\textbf{Dịch vụ} & \textbf{Trách nhiệm kỹ thuật} \\
\midrule
\texttt{web} & Giao diện Next.js cho upload audio, xem lịch sử meeting, review transcript/task và gửi feedback. \\
\texttt{api} & FastAPI REST API, tạo meeting job, đọc/ghi PostgreSQL, xuất Prometheus metrics và xử lý admin endpoints. \\
\texttt{worker} & Celery worker chạy pipeline nặng: ingest, preprocess, ASR, diarization, LLM extraction, guardrails và persistence. \\
\texttt{postgres} & Lưu meetings, transcript turns, tasks, workers, feedback corrections, artifacts và calendar events. \\
\texttt{redis} & Celery broker/result backend và cache cho một số bước inference. \\
\texttt{ollama} & Runtime phục vụ Qwen2.5-3B và embedding model local. \\
\texttt{prometheus}, \texttt{grafana} & Thu thập và hiển thị metrics vận hành. \\
\texttt{beat}, \texttt{trainer}, \texttt{mlflow} & Profile MLOps cho retraining schedule, training worker và experiment tracking. \\
\bottomrule
\end{tabularx}
\end{table}

\section{Tương tác người dùng và định dạng vào ra}

Luồng người dùng gồm bốn bước: upload audio qua web UI, theo dõi trạng thái job, review transcript/action items, và gửi feedback correction. Đầu vào chính là audio file và optional roster. Đầu ra chính là summary, transcript turns và action items có các trường \texttt{description}, \texttt{assignee}, \texttt{due\_date}, \texttt{priority}, \texttt{status} và \texttt{source\_turn\_ids}.

Backend lưu trạng thái job trong PostgreSQL thay vì giữ trong bộ nhớ worker, nên frontend có thể polling \texttt{GET /meetings/\{id\}} để theo dõi job dài. Đường dẫn \texttt{GET /metrics} xuất Prometheus metrics; các đường dẫn quản trị hỗ trợ kiểm tra retraining và trạng thái MLOps.

Bảng \ref{tab:api-io} trình bày các API chính, định dạng đầu vào/đầu ra và phần hệ thống chịu trách nhiệm xử lý.


\begin{table}[H]
\centering
\footnotesize
\caption{Các API chính và định dạng dữ liệu}
\label{tab:api-io}
\begin{tabularx}{\textwidth}{@{}>{\raggedright\arraybackslash}p{2.9cm}>{\raggedright\arraybackslash}p{3.9cm}X@{}}
\toprule
\textbf{Method} & \textbf{Đường dẫn} & \textbf{Vai trò và dữ liệu} \\
\midrule
\texttt{POST} & \apipath{/meetings} & Nhận audio file và \texttt{roster\_json}; trả \texttt{meeting\_id} để theo dõi job bất đồng bộ. \\
\texttt{GET} & \apipath{/meetings/{id}} & Trả trạng thái meeting, transcript turns, summary, action items, unresolved items, human review items và run metrics. \\
\texttt{POST} & \apipath{/meetings/{id}/feedback} & Nhận correction cho task: description, assignee, due date, false positive hoặc missing task. \\
\texttt{POST} & \apipath{/meetings/{id}/participants/{speaker}/resolve} & Map speaker diarization label sang tên người tham gia hoặc worker trong roster. \\
\scriptsize\texttt{GET/POST/PUT/DELETE}\normalsize & \apipath{/workers} & Quản lý worker roster gồm name, email, role, aliases và skills. \\
\texttt{GET} & \apipath{/metrics}; \apipath{/metrics/business} & Xuất metrics vận hành và business metrics cho Prometheus/Grafana. \\
\texttt{POST/GET} & \apipath{/admin/retrain}; \apipath{/admin/retrain/state} & Kiểm tra hoặc kích hoạt retraining pipeline và xem trạng thái retraining. \\
\bottomrule
\end{tabularx}
\end{table}


\section{Latency, khả năng mở rộng và versioning mô hình}

Pipeline được thiết kế bất đồng bộ vì ASR, diarization và LLM inference có độ trễ cao hơn request HTTP thông thường. Redis/Celery cho phép tách API nhận request khỏi worker xử lý nặng. Model version được ghi vào metadata meeting và promotion flow dùng manifest để kiểm soát việc thay đổi model. Cách triển khai này phù hợp cho demo local và có đường mở rộng sang queue scaling, GPU trainer và model promotion có kiểm soát.
